\section{Introduction}
Thin films are excellent systems to research electroceramic materials. Modern deposition
methods enable the engineering of complex thin films, such as multilayers, nanocomposites
or layers with a graded composition. This results in a variety of applications, including
microelectronics, optics and information storage. In most cases, the chemical
and physical properties between the thin film and its bulk material are
different. These differences include, but are not limited to: geometric differences,
grain sizes, preferential orientation, crystallinity differences, strain / stress,
compositional gradients,

Most of them originate from the growth conditions which can differ for each
deposition process. It is therefore necessary to use tools to routinely characterize
newly deposited thin films. A useful and versatile technique for thin films is x-ray
diffraction. It can provide information about crystalline phase and lattice parameters,
grain size, density of dislocations, residual stress/strain, film composition,

XRD is a non-destructive measurement technique that yields precise and accurate results
if carried out properly. However, it has a poor spatial resolution, probing a
few mm² to cm².